\section{\texorpdfstring{\(J\)}{J}-normalization at discrepancy one}
\label{sec:normalization}

Let \(\phi\colon Y\to Z\) be the flag bundle of \(V\oplus V'\), so that
\(\phi^{*}V\) and \(\phi^{*}V'\) carry filtrations whose successive quotients
are line bundles, and \(\phi^{*}\colon H^{*}(Z)\to H^{*}(Y)\) is injective;
this is the splitting principle for extremal quantum cohomology,
\cite[Corollary~4.3]{SS}.  The Chern roots \(\rho_{i}\) and \(\sigma_{j}\)
exist individually on \(Y\), and we continue to write \(V\), \(V'\),
\(\rho_{i}\), \(\sigma_{j}\) for the pullbacks.  Write
\[
  A=H^{*}(\PP_{Y}(V))
   \cong H^{*}(Y)[H]\big/\left\langle
     \textstyle\prod_{i=1}^{r}(\rho_i+H)\right\rangle
\]
for the coefficient ring of \eqref{eq:ss35} after this pullback, and let
\(R=A[\ln q][[t]]\), the coordinates of \(t\) being formal variables; the
\(\ln q\) is discussed before Lemma~\ref{lem:jslice}.  Each factor of the
denominator of \eqref{eq:ss35} has the form \(mz+a\) with \(m\ge1\) and
\(a=\rho_{i}+H\in A\), and we expand
\[
  \frac1{mz+a}
  =\frac1{mz}\sum_{l\ge0}\Bigl(\frac{-a}{mz}\Bigr)^{l}\in z^{-1}A[[z^{-1}]].
\]
The sum is finite because elements of \(A\) of positive degree are nilpotent,
but only the formal expansion is used.  The prefactor
\(z\,e^{t/z}q^{c_{1}(T)/z}=z\,e^{\bar{\mathbf t}/z}\) lies in
\(z\one+R[[z^{-1}]]\) and is an infinite series in \(z^{-1}\), since \(t\) has
formal coordinates.  So the right-hand side of \eqref{eq:ss35} lies in
\(zR[[z^{-1}]][[q]]\).
By the \emph{\(z\)-order} of a nonzero element of \(zR[[z^{-1}]]\) we mean the
largest \(n\) with a nonzero \(z^{n}\) coefficient.  A \(z\)-order bound proved after the pullback
is a coefficientwise statement about classes in the image of an injective
map, so it descends to the original cohomology ring; the bound below is in
any case stated in terms of \(r\), \(s\) and \(d\) alone.

\begin{lemma}
  \label{lem:order}
  For every \(d\ge1\), the degree-\(d\) summand of \eqref{eq:ss35} has
  \(z\)-order at most
  \[
    1+s(d-1)-rd=1-s-\nu d .
  \]
\end{lemma}

\begin{proof}
  The numerator \(\prod_{j=1}^{s}\prod_{m=0}^{d-1}(\sigma_{j}-H-mz)\) has
  \(sd\) factors.  The \(s\) factors with \(m=0\) are \(\sigma_{j}-H\) and
  contain no \(z\); each of the remaining \(s(d-1)\) factors is linear in
  \(z\).  So the numerator is a polynomial in \(z\) over \(A\) of degree at
  most \(s(d-1)\).  The denominator \(\prod_{i=1}^{r}\prod_{m=1}^{d}(\rho_{i}
  +H+mz)\) has \(rd\) factors, all with \(m\ge1\), so by the expansion above
  its inverse lies in \(z^{-rd}A[[z^{-1}]]\).  The prefactor
  \(z\,e^{\bar{\mathbf t}/z}\) is \(z\) times an element of
  \(\one+z^{-1}R[[z^{-1}]]\); the factor \(z\) adds one to the order, and the
  second factor does not raise it.
\end{proof}

The two hypotheses of Theorem~\ref{thm:normalization} enter here and nowhere
else.

\begin{corollary}
  \label{cor:normalized}
  If \(\nu=1\) and \(s\ge1\), then every summand of \eqref{eq:ss35} with
  \(d\ge1\) has \(z\)-order at most \(1-s-d\le-1\), and hence
  \[
    z\,e^{t/z}q^{c_{1}(T)/z}
    \sum_{d\ge0}q^{(r-s)d}\,
    \frac{\prod_{j=1}^{s}\prod_{m=0}^{d-1}(\sigma_{j}-H-mz)}
         {\prod_{i=1}^{r}\prod_{m=1}^{d}(\rho_{i}+H+mz)}
    =z\one+\bar{\mathbf t}+O(z^{-1}),
  \]
  with \(\bar{\mathbf t}=t+\ln(q)c_{1}(T)\) as in \eqref{eq:barparam}.
\end{corollary}

\begin{proof}
  The bound of Lemma~\ref{lem:order} is \(1-s-d\), which is at most \(-1\)
  for \(d\ge1\) and \(s\ge1\).  The \(d=0\) summand is
  \(z\,e^{t/z}q^{c_{1}(T)/z}=z\,e^{\bar{\mathbf t}/z}
  =z\one+\bar{\mathbf t}+O(z^{-1})\).
\end{proof}

\begin{remark}
  \label{rem:boundary}
  Both hypotheses are sharp.  For \(\nu=1\) and \(s=0\), so \((r,s)=(1,0)\),
  the bound of Lemma~\ref{lem:order} is \(1-d\) and is attained by the formal
  \(d=1\) term of the hypergeometric expression.  Here \(\rho_{1}+H=0\) in
  \(A\), so that term is exactly
  \[
    z\cdot\frac{q}{\rho_{1}+H+z}=q\one .
  \]
  Its \(z^{0}\) coefficient is \(q\one\), so this formal expression is not
  \(J\)-normalized.  Geometrically \(\PP(V)=Z\) has point fibres and no line
  class \([L]\), so \((1,0)\) is a degenerate endpoint of the extremal-line
  setup, not a further standard-flip case repaired here.  For \(\nu=0\) the bound is the
  \(d\)-independent value \(1-s\) and the sum is not \(q\)-adically
  convergent, consistently with the standing assumption \(r>s\) in \cite{SS}.
\end{remark}

To convert the normalization into an identification we use the standard
characterization of the \(J\)-function as a slice of Givental's cone.  The
space \(T\) is not compact, so \(H^{*}(T)\) carries no nondegenerate
Poincar\'e pairing and the symplectic formalism is available only
equivariantly.  Let \(\C^{\times}\) act trivially on \(\PP(V)\) and by inverse
scaling on the fibres of \(V'(-1)=\pi^{*}V'\otimes\mathcal O_{\PP(V)}(-1)\),
and write \(\lambda\) for the generator of \(H^{2}_{\C^{\times}}(\mathrm{pt})\)
given by the standard character.  The equivariant Chern roots of \(V'(-1)\)
are then \(\sigma_{j}-H-\lambda\), so
\(e_{\lambda}(V'(-1))=\prod_{j=1}^{s}(-\lambda+\sigma_{j}-H)\); this is the
only sign convention for \(\lambda\) used below.  Following
\cite{CoatesGivental} and \cite{CCIT}, the local theory of \(T\) is
the genus-zero theory of the compact space \(\PP(V)\) twisted by the inverse
\(\C^{\times}\)-equivariant Euler class of \(V'(-1)\); its pairing
\((a,b)_{\lambda}=\int_{\PP(V)}ab\,e_{\lambda}(V'(-1))^{-1}\) is
nondegenerate over \(\C(\lambda)\).  Let \(\mathcal L_{T}^{\lambda}\) and
\(J^{T}_{\lambda}\) denote its overruled Lagrangian cone and its
\(J\)-function, and \(\mathcal L_{\PP(V)}\) the cone of \(\PP(V)\), all in
the extremal specialization \eqref{eq:barparam}.  The extremal
\(J\)-function \(J^{T}\) of \cite{SS} is recovered at \(\lambda=0\) by
Lemma~\ref{lem:limit} below.

The parameter \(\bar{\mathbf t}=t+\ln(q)c_{1}(T)\) is not in the maximal
ideal of the formal parameters, and we fix its meaning by the divisor
equation.  For a theory \(\mathcal T\) (untwisted, or twisted as above),
with \(J\)-function \(J^{\mathcal T}=\sum_{d}Q^{d}J^{\mathcal T}_{d}\), a
divisor class \(D\) and \(t\) formal, we set
\begin{equation}
  \label{eq:lnq}
  J^{\mathcal T}(t+\ln(q)D,z)
  :=e^{\ln(q)D/z}\sum_{d}Q^{d}q^{\langle D,d\rangle}J^{\mathcal T}_{d}(t,z),
\end{equation}
which is what the divisor equation gives whenever both sides are defined; the
divisor equation holds for the twisted theory as for the untwisted one
\cite{CoatesGivental}.  With \(D=c_{1}(T)\) and the extremal specialization,
the degree-\(k[L]\) summand carries \(q^{(r-s)k}\), and the factor
\(e^{\ln(q)c_{1}(T)/z}=q^{c_{1}(T)/z}\) is the prefactor of
\eqref{eq:ss35}.  The weight \(q^{(r-s)k}\) is positive because \(r>s\), so
the series converges \(q\)-adically.  The same divisor-equation symmetry,
multiplication by \(e^{\ln(q)D/z}\) together with \(Q^{d}\mapsto
Q^{d}q^{\langle D,d\rangle}\), preserves each cone below, so a statement at
the parameter \(\bar{\mathbf t}\) is the image of the same statement at the
formal parameter \(t\).

\begin{lemma}[\(J\)-slice uniqueness; {\cite{Givental}}, {\cite{CoatesGivental}}]
  \label{lem:jslice}
  Let \(\mathcal L\) be \(\mathcal L_{\PP(V)}\) or
  \(\mathcal L_{T}^{\lambda}\), with \(J\)-function \(J\), and let
  \(f\in\mathcal L\) be of the form \(f=z\one+\tau+O(z^{-1})\) for
  some \(\tau\) in the parameter space.  Then \(f=J(\tau,z)\).
\end{lemma}

For \(\tau=\bar{\mathbf t}\) the lemma is applied in the form fixed by
\eqref{eq:lnq}: strip the factor \(q^{c_{1}(T)/z}\) from \(f\), apply
uniqueness at the formal parameter \(t\), and restore the factor.

This is the mechanism behind every mirror theorem of Givental type: a point
of the cone with a prescribed \(z^{0}\) coefficient is the \(J\)-function at
that coefficient, and the mirror map is precisely the difference between
the prescribed \(z^{0}\) coefficient and the parameter one started with.

The second ingredient is that \eqref{eq:ss35}, in its \(\lambda\)-twisted
form, lies on the cone at all.  This is asserted in
\cite[Remark~4.5(3)]{SS}, the same remark whose other assertion
Corollary~\ref{cor:normalized} contradicts, so we prove it instead of quoting
it.  We begin with split bundles, where the two sources the authors name do
the work, with the inverse-Euler twist taken from \cite{CCIT}.

\begin{lemma}
  \label{lem:cone}
  Let \(r>s\ge0\) with \(r\ge2\), and suppose \(V=\bigoplus_{i=1}^{r}L_{i}\)
  and \(V'=\bigoplus_{j=1}^{s}M_{j}\).  Then the series
  \begin{equation}
    \label{eq:Ilambda}
    I_{\lambda}
    =z\,e^{t/z}q^{c_{1}(T)/z}
    \sum_{d\ge0}q^{(r-s)d}\,
    \frac{\prod_{j=1}^{s}\prod_{m=0}^{d-1}(-\lambda+\sigma_{j}-H-mz)}
         {\prod_{i=1}^{r}\prod_{m=1}^{d}(\rho_{i}+H+mz)}
  \end{equation}
  lies on \(\mathcal L_{T}^{\lambda}\).  At \(\lambda=0\) it is the
  right-hand side of \eqref{eq:ss35}.
\end{lemma}

\begin{proof}
  The projective bundle \(\PP(V)\to Z\) is a toric fibration with fibre
  \(\PP^{r-1}\).  Brown's theorem \cite[Theorem~1]{Brown} states that the
  hypergeometric modification of the \(J\)-function of the base lies on the
  overruled cone of the total space, and \cite[Corollary~1]{Brown} is the
  projective-fibration case.  Setting the Novikov variables of \(Z\) to zero
  and keeping the fibrewise variable is the extremal specialization
  \cite[(33)]{SS}, a homomorphism of Novikov rings by
  \cite[Proposition~4.1]{SS}; it leaves \cite[(34)]{SS}, which therefore
  lies on the extremal cone \(\mathcal L_{\PP(V)}\) of \(\PP(V)\).  Brown's
  hypotheses concern the fibration
  \(\PP(V)\to Z\) and the Chern roots \(\rho_{1},\dots,\rho_{r}\); no second
  bundle occurs in them, so they cannot constrain \(r-s\).

  Now the twist.  A stable map \(f\colon C\to\PP(V)\) of class \(k[L]\) with
  \(k\ge1\) has image in a fibre of \(\PP(V)\to Z\), by the argument of
  \cite[Lemma~2.6]{SS} in the fibre-bundle form used in the proof of
  \cite[Theorem~4.2]{SS} (the lemma itself is stated for curves in \(X\)).
  There \(V'(-1)=\pi^{*}V'\otimes\mathcal O_{\PP(V)}(-1)\) pulls back to a
  sum of \(s\) line bundles of degree \(-k_{i}\) on each component
  \(C_{i}\), where \(k_{i}\ge0\) is the degree of \(f|_{C_{i}}\) and
  \(\sum_{i}k_{i}=k\).  A section vanishes on every component with
  \(k_{i}>0\), and, the dual graph being a connected tree with at least one
  such component, on the contracted components as well; so
  \(H^{0}(C,f^{*}V'(-1))=0\) and
  \(V'(-1)\) is concave along every extremal class, while its dual
  \(W:=(V'(-1))^{\vee}=\pi^{*}(V')^{\vee}\otimes\mathcal O_{\PP(V)}(1)\) is
  convex there.  The twisted theory is the theory of \(\PP(V)\) twisted by the
  invertible multiplicative class \(c=e_{\lambda}^{-1}\) and the split bundle
  \(F=V'(-1)=\bigoplus_{j}\pi^{*}M_{j}\otimes\mathcal O_{\PP(V)}(-1)\).  For
  such a pair, \cite[Theorem~4.6]{CCIT} proves that the hypergeometric
  modification \(I^{\mathrm{tw}}=\sum_{D}Q^{D}M_{D}(z)J^{\PP(V)}_{D}\) of the
  \(J\)-function of \(\PP(V)\) lies on the twisted cone; it holds for an
  arbitrary invertible multiplicative class and a direct sum of line bundles,
  and the inverse equivariant Euler class of a concave bundle is the case
  \cite[Example~A]{CCIT}.  For the Euler class of a convex bundle it is
  \cite[Theorem~2]{CoatesGivental}, and the inverse-Euler case is the
  Serre-dual partner discussed in ``On Serre duality'' of
  \cite{CoatesGivental}.  The theorem is a statement over the full Novikov
  ring of \(\PP(V)\), with \(J^{\PP(V)}_{D}\) the degree-\(D\) coefficient of
  \(J^{\PP(V)}\) for every effective class \(D\).

  The modification factor of a line bundle with first Chern class \(x\) in
  degree \(D\) is
  \(M_{D}=\prod_{m\le\langle x,D\rangle}c(x+mz)\big/\prod_{m\le0}c(x+mz)\)
  \cite[Section~4]{CCIT}.  Here the equivariant roots of \(F\) are
  \(x_{j}=\sigma_{j}-H-\lambda\), and \(\langle x_{j},d[L]\rangle=-d\), so with
  \(c(x)=x^{-1}\) the factor contributed by the \(j\)-th root is
  \(\prod_{l=-d+1}^{0}(-\lambda+\sigma_{j}-H+lz)\), which is empty for
  \(d=0\).  Writing \(m=-l\) turns this into
  \(\prod_{m=0}^{d-1}(-\lambda+\sigma_{j}-H-mz)\), the numerator of
  \eqref{eq:Ilambda}.

  We now pass to the extremal specialization.  Setting the Novikov variables
  of classes outside \(\mathbb N[L]\) to zero is a ring homomorphism by
  \cite[Proposition~4.1]{SS}, since those classes form a monoid ideal of the
  effective cone.  The twisted cone is defined by generating functions taken
  coefficientwise in the Novikov degree, so the homomorphism carries
  \(I^{\mathrm{tw}}\) to a point of the extremal twisted cone
  \(\mathcal L^{\lambda}_{T}\), and only the terms \(D=d[L]\) survive.  It
  remains to identify the coefficients \(J^{\PP(V)}_{d[L]}\) with those of
  \cite[(34)]{SS}.  For \(r\ge2\) the count of Lemma~\ref{lem:order},
  applied with \(s=0\), gives the degree-\(d\) summand of \cite[(34)]{SS}
  \(z\)-order at most \(1-rd\le-1\) for every \(d\ge1\); so
  \cite[(34)]{SS} is \(J\)-normalized and, lying on \(\mathcal L_{\PP(V)}\),
  equals the extremal specialization of \(J^{\PP(V)}\) by
  Lemma~\ref{lem:jslice}.  Reparametrizing by \eqref{eq:lnq} with
  \(D=c_{1}(T)\), as in \cite[(33)]{SS}, the class \(d[L]\) carries the
  weight \(q^{\langle c_{1}(T),\,d[L]\rangle}=q^{(r-s)d}\) and the prefactor
  becomes \(q^{c_{1}(T)/z}\); this gives \eqref{eq:Ilambda} as a point of
  \(\mathcal L_{T}^{\lambda}\).  The modification factor carries no \(q\);
  the exponent changes only because the reparametrization is by
  \(c_{1}(T)=c_{1}(\PP(V))+c_{1}(V'(-1))\).  The \(d=0\) summand is untouched
  by the modification and remains the classical term
  \(z\,e^{\bar{\mathbf t}/z}\).  The modification involves the Chern roots of
  \(V'(-1)\) and nothing else; here too \(r\) and \(s\) enter only as numbers
  of factors.
\end{proof}

The cone \(\mathcal L_{T}^{\lambda}\) is defined over \(\C(\lambda)\), and its
pairing degenerates at \(\lambda=0\), so the comparison with \(J^{T}\) must be
made after \(J\)-slice uniqueness has been applied, not before.  The next
lemma makes that comparison.

\begin{lemma}
  \label{lem:limit}
  Let \(r\ge2\).  In every degree \(d=k[L]\) with \(k\ge1\), the
  coefficients of \(J^{T}_{\lambda}\) are polynomial in \(\lambda\), and
  their values at \(\lambda=0\) are the coefficients of the extremal
  \(J\)-function \(J^{T}\) of \cite{SS}.  In degree zero both are
  \(z\,e^{\bar{\mathbf t}/z}\).
\end{lemma}

\begin{proof}
  Write \(E=V'(-1)\) and \(E_{0,n+1,d}=R^{0}\pi_{\mathcal C*}f^{*}E
  -R^{1}\pi_{\mathcal C*}f^{*}E\).  For \(k\ge1\) the vanishing
  \(H^{0}(C,f^{*}E)=0\) shown in the proof of Lemma~\ref{lem:cone} holds for
  every stable map of class \(k[L]\) with any number of marked points, so the
  twisting class is
  \(e_{\lambda}^{-1}(E_{0,n+1,d})=e_{\lambda}(R^{1}\pi_{\mathcal C*}f^{*}E)\),
  a polynomial in \(\lambda\).  The basis dual to \(\{\phi_{a}\}\) for
  \((\ ,\ )_{\lambda}\) is \(\{e_{\lambda}(E)\phi^{a}\}\), with
  \(\{\phi^{a}\}\) dual for the pairing of \(\PP(V)\).  So the
  degree-\(d\) part of \(J^{T}_{\lambda}\) is
  \(e_{\lambda}(E)\cup\mathrm{ev}_{n+1,*}\bigl(e_{\lambda}(R^{1}\pi_{\mathcal
  C*}f^{*}E)\cap\cdots\bigr)\), summed over \(n\), with the descendant
  insertions of \(J\) in place of the dots; this is polynomial in \(\lambda\).
  At \(\lambda=0\) it becomes \(e(E)\cup\mathrm{ev}_{n+1,*}\bigl(e(R^{1}\pi_{\mathcal
  C*}f^{*}E)\cap\cdots\bigr)\).  The Euler class of the obstruction bundle is
  the virtual class of
  \(\overline M_{0,n+1}(T,d)=\overline M_{0,n+1}(\PP(V),d)\): this is
  \cite[Corollary~3.2]{SS}, stated there for \(X\), with \(X\) replaced by
  \(T\) as in the proof of \cite[Theorem~3.3]{SS}.  Moreover
  \(e(E)\cup-\) is \(j_{T}^{*}j_{T*}\), which identifies a pushforward to
  \(\PP(V)\) with the pushforward to \(T\) along the proper evaluation map,
  again as in that proof.  The result is the degree-\(d\) part
  of \(J^{T}\).  This is the non-equivariant specialization in positive degree
  recorded in the discussion ``On Serre duality'' of \cite{CoatesGivental}.  In
  degree zero the twist \(e_{\lambda}(E)^{-1}\) and the factor
  \(e_{\lambda}(E)\) of the dual basis cancel, leaving the classical term.
\end{proof}

The passage from split bundles to general ones is where the printed
attribution of \cite[Remark~4.5(3)]{SS} cannot be followed at \(\nu=1\), so
we spell it out.  What it transports is the identity between
\eqref{eq:ss35} and \(J^{T}\), not membership in an infinite-dimensional
cone.

\begin{lemma}
  \label{lem:split-reduction}
  Let \(r>s\ge0\) with \(r\ge2\).  Suppose the right-hand side of
  \eqref{eq:ss35} equals the extremal \(J\)-function for every split local
  model with the same \((r,s)\) over every smooth projective base.  Then it
  does so for every local model with that \((r,s)\).
\end{lemma}

The hypothesis \(r\ge2\) only guarantees that the line class \([L]\) exists;
it is not used otherwise in the proof.

\begin{proof}
  Let \(\phi\colon Y\to Z\) be the flag bundle of \(V\oplus V'\) and
  \(\tilde\phi\colon T_{Y}=Y\times_{Z}T\to T\) the induced map.  The
  Cartesian square identifies the relative obstruction theories and hence
  the virtual classes of \(\overline M_{0,n}(T_{Y},d)\) and
  \(\overline M_{0,n}(T,d)\) for \(d\) in the span of \([L]\); this is the
  content of \cite[Theorem~4.2]{SS} and, for the fundamental solution, of
  \cite[Lemma~9.5]{SS}.  In positive degree \(J^{T}\) is the pushforward
  along the proper evaluation map of descendant classes on these moduli
  spaces, so base change in the Cartesian square gives, at every formal
  parameter \(\tau\) of \(T\),
  \(\tilde\phi^{*}J^{T}(\tau)=J^{T_{Y}}(\tilde\phi^{*}\tau)\).

  The parameters do not match exactly.  Write \(\Theta\) for tangent bundles
  and \(\varpi\colon T_{Y}\to Y\) for the projection.  The relative tangent
  bundle of \(\tilde\phi\) is the base change \(\varpi^{*}\Theta_{Y/Z}\), so
  \(0\to\varpi^{*}\Theta_{Y/Z}\to\Theta_{T_{Y}}\to\tilde\phi^{*}\Theta_{T}\to0\)
  is exact and
  \[
    c_{1}(T_{Y})=\tilde\phi^{*}c_{1}(T)+D_{Y},\qquad
    D_{Y}:=\varpi^{*}c_{1}(\Theta_{Y/Z}),
  \]
  so the parameter \(\bar{\mathbf t}_{Y}=\tilde\phi^{*}t+\ln(q)c_{1}(T_{Y})\)
  of \(T_{Y}\) is \(\tilde\phi^{*}\bar{\mathbf t}+\ln(q)D_{Y}\).  The class
  \(D_{Y}\) is pulled back from \(Y\), and a line in a fibre of
  \(\PP(V_{Y})\to Y\) maps to a point of \(Y\), so
  \(\langle D_{Y},[L]\rangle=0\).  By the divisor equation \eqref{eq:lnq},
  \[
    J^{T_{Y}}(\bar{\mathbf t}_{Y})
    =e^{\ln(q)D_{Y}/z}\,J^{T_{Y}}(\tilde\phi^{*}\bar{\mathbf t})
    =e^{\ln(q)D_{Y}/z}\,\tilde\phi^{*}J^{T}(\bar{\mathbf t}).
  \]
  On the other side, both numerator and denominator of \eqref{eq:ss35} are
  symmetric in the respective Chern roots; \cite[Remark~4.5(2)]{SS} says this
  for \cite[(34)]{SS}, and the same holds for \eqref{eq:ss35}.  The
  summands for \(T_{Y}\) are therefore the pullbacks of those for \(T\), while
  the prefactor is
  \(q^{c_{1}(T_{Y})/z}=e^{\ln(q)D_{Y}/z}\,\tilde\phi^{*}q^{c_{1}(T)/z}\).
  So the expression \eqref{eq:ss35} for \(T_{Y}\) is
  \(e^{\ln(q)D_{Y}/z}\) times the pullback of the expression for \(T\).  The
  same invertible factor multiplies both sides, and \(\tilde\phi^{*}\) is
  injective by Leray--Hirsch \cite[Corollary~4.3]{SS}; so it suffices to prove
  the identity on \(T_{Y}\), where \(V\) and \(V'\) are filtered with
  line-bundle quotients.  The base \(Y\) is again smooth projective.

  It remains to replace a filtered bundle by its associated graded.  Take one
  step \(0\to A\to V\to Q\to0\) of the filtration; in this paragraph \(A\),
  \(Q\) and \(t\) are the sub- and quotient bundles and the coordinate on
  \(\mathbf A^{1}\), not the coefficient ring above, the Novikov variable of
  \eqref{eq:jnorm} or the parameter of \eqref{eq:barparam}.  Let
  \(d\colon V\to Q\) be the quotient map and, on \(Z\times\mathbf A^{1}\),
  \[
    g\colon V\oplus Q\to Q,\qquad (v,w)\mapsto d(v)-t\,w .
  \]
  It is surjective, so \(\ker(g)\) is a bundle of rank \(r\), with
  \(\ker(g)_{1}=\{(v,d(v))\}\cong V\) and \(\ker(g)_{0}=A\oplus Q\).  This
  is the construction of \cite[Lemma~9.4]{SS}, where the domain is printed as
  \(A\oplus Q\); with that domain \(d\) vanishes and \(\ker(g)_{1}\) would be
  \(A\).  We obtain a smooth projective family
  \(\PP(\ker g)\to Z\times\mathbf A^{1}\) carrying the bundle \(V'(-1)\).
  For \(k\ge1\), the argument of
  \cite[Proposition~3.1]{SS}, stated there for \(X\), applied to the projective bundle
  \(\PP(\ker g)\to Z\times\mathbf A^{1}\), makes the moduli spaces
  \(\overline M_{0,n}(\PP(\ker g)_{t},k[L])\) a smooth proper family
  over \(\mathbf A^{1}\), and concavity makes
  \(R^{1}\pi_{\mathcal C*}f^{*}V'(-1)\), for \(\pi_{\mathcal C}\) and \(f\)
  the universal curve and universal map, a relative locally free obstruction
  bundle of rank \(s(k-1)\).  Its Euler class, capped with the relative
  fundamental class, gives a relative cycle whose fibrewise restrictions
  are the extremal virtual classes, by the argument of
  \cite[Corollary~3.2]{SS} on each fibre.  The relative
  evaluation maps and \(\psi\)-classes define relative descendant cycles, and
  pushing these to \(\mathbf A^{1}\) gives classes in
  \(H^{0}(\mathbf A^{1})\), hence constants.  So the extremal genus-zero
  descendant invariants with insertions restricted from
  \(H^{*}(\PP(\ker g))\) agree at \(t=0\) and \(t=1\).  These determine the
  extremal \(J\)-functions.  By Leray--Hirsch, \(H^{*}(\PP(\ker g)_{t})\) is
  spanned by restrictions of classes on \(\PP(\ker g)\), namely powers of
  \(H\) times classes from \(Z\), and the factor \(e(V'(-1))\) of the
  pushforward form of \(J^{T}\) in Lemma~\ref{lem:limit} is the restriction
  of a family class as well.  The degree-\(k[L]\) coefficient of \(J^{T}\) is
  therefore fixed by its Poincar\'e pairings on the compact space
  \(\PP(\ker g)_{t}\) with a basis of family classes, and those pairings are
  the constant invariants just obtained.  So the extremal \(J\)-functions
  agree under the isomorphism of cohomology rings of \cite[Lemma~9.4]{SS}.  The right-hand side of \eqref{eq:ss35} is likewise
  unchanged, because \(V\) and \(A\oplus Q\) have the same Chern roots.
  Repeating the construction for the filtration of \(V'\), and iterating both
  filtrations as in \cite[Lemma~9.6]{SS}, reaches the split
  associated-graded local model.
\end{proof}

\begin{proposition}
  \label{prop:IJ}
  Let \(r=s+1\) with \(s\ge1\).  Then the right-hand side of
  \eqref{eq:ss35} equals the extremal \(J\)-function \(J^{T}(q,t,z)\) of the
  local model, in the variables and completion of \cite[Theorem~4.4]{SS}.
\end{proposition}

\begin{proof}
  Assume first that \(V\) and \(V'\) split.  By Lemma~\ref{lem:cone} the
  series \eqref{eq:Ilambda} is a point of \(\mathcal L_{T}^{\lambda}\).  The
  count of Lemma~\ref{lem:order} applies to it unchanged: \(\lambda\) enters
  only the numerator factors, which keep their degree in \(z\), and the
  \(s\) factors with \(m=0\), now \(-\lambda+\sigma_{j}-H\), are still free
  of \(z\).  So, as in Corollary~\ref{cor:normalized},
  \(I_{\lambda}=z\one+\bar{\mathbf t}+O(z^{-1})\), and by uniqueness of the
  \(J\)-slice, Lemma~\ref{lem:jslice} with \(\tau=\bar{\mathbf t}\),
  \(I_{\lambda}=J^{T}_{\lambda}(q,t,z)\).  Both sides are polynomial in
  \(\lambda\) in each degree, \(I_{\lambda}\) visibly and \(J^{T}_{\lambda}\)
  by Lemma~\ref{lem:limit}.  Setting \(\lambda=0\) gives \eqref{eq:ss35} on
  the left and \(J^{T}(q,t,z)\) on the right, by Lemmas~\ref{lem:cone}
  and~\ref{lem:limit}.  The general case follows from
  Lemma~\ref{lem:split-reduction}.
\end{proof}

The Shen--Shoemaker argument after their Section~4 uses \eqref{eq:ss35} also in
equivariant cohomology.  Their Lemma~7.4 and Proposition~9.1 use it for the
torus \(\mathbb T=(\C^{\times})^{r+s}\) scaling the summands of split bundles
\cite[Sections~7.2 and~9.1]{SS}.  The proof of their Lemma~9.4 uses it for the
torus \(S=(\C^{\times})^{2}\) scaling \(V\) and \(V'\).  The equivariant form
follows from Proposition~\ref{prop:IJ} by approximation, because extremal
curves never leave the fibres over \(Z\).

\begin{corollary}
  \label{cor:equivariant}
  Let \(r=s+1\) with \(s\ge1\), and let \(G\) be a torus acting trivially on
  \(Z\) and by vector-bundle automorphisms on \(V\) and \(V'\).  Then
  \eqref{eq:ss35}, with \(\rho_{i}\), \(\sigma_{j}\), \(H\) and \(c_{1}(T)\)
  read as \(G\)-equivariant classes and \(t\in H^{*}_{G}(Z)\), equals the
  \(G\)-equivariant extremal \(J\)-function of \(T\).  Here the latter is
  defined in degree \(k[L]\), \(k\ge1\), by the proper-pushforward formula
  \(\mathrm{ev}_{*}\bigl(e_{G}(R^{1}\pi_{\mathcal C*}f^{*}V'(-1))\cap\cdots\bigr)\)
  along \(\mathrm{ev}\colon\overline M_{0,n+1}(T,k[L])\to T\), and in degree
  zero by \(z\,e^{\bar{\mathbf t}/z}\).  When the \(G\)-fixed locus of \(T\)
  is compact, this agrees with the definition through the localized
  equivariant pairing.  This applies to \(G=S\), and to \(G=\mathbb T\) when
  \(V\) and \(V'\) split.
\end{corollary}

The pushforward definition is needed in general: if \(G\) acts on \(V\) and
\(V'\) with weights that make the weight of \(V'(-1)\) trivial, the fixed
locus is all of \(T\) and the localized pairing is undefined.  For \(S\) and
\(\mathbb T\) the fixed locus lies in the zero section \(\PP(V)\), which is
compact.

\begin{proof}
  Write \(G\cong(\C^{\times})^{g}\), \(E_{N}=(\C^{N}\setminus0)^{g}\) and
  \(B_{N}=E_{N}/G=(\PP^{N-1})^{g}\).  Then \(V_{N}=V\times_{G}E_{N}\) and
  \(V'_{N}=V'\times_{G}E_{N}\) are vector bundles of ranks \(r\) and \(s\) on
  the smooth projective variety \(Z\times B_{N}\), and
  \(T_{N}=T\times_{G}E_{N}\) is their local model.  Restriction
  \(H^{k}_{G}(T)\to H^{k}(T_{N})\) is an isomorphism for \(k<2N-1\).  It sends
  the equivariant classes \(\rho_{i}\), \(\sigma_{j}\), \(H\) to the
  corresponding classes of \(V_{N}\), \(V'_{N}\) and \(\mathcal
  O_{\PP(V_{N})}(1)\), and \(c^{G}_{1}(T)\) to
  \(c_{1}(T_{N})-c_{1}(B_{N})\).

  By the argument of \cite[Lemma~2.6]{SS} in the fibre-bundle form of the
  proof of \cite[Theorem~4.2]{SS}, applied over the projective base
  \(Z\times B_{N}\), stable maps of class \(k[L]\) to \(\PP(V_{N})\) lie in
  fibres over \(Z\times B_{N}\); by the vanishing of sections in the proof of
  Lemma~\ref{lem:cone}, stable maps of that class to \(T_{N}\) lie on the
  zero section.  So
  \(\overline M_{0,n}(T_{N},k[L])=\overline M_{0,n}(T,k[L])\times_{G}E_{N}\),
  and the obstruction bundle (that of \cite[Corollary~3.2]{SS}, with \(X\)
  replaced by \(T\)) descends.  These
  moduli spaces are smooth and their virtual classes are Euler classes.  The
  equivariant extremal descendant classes of \(T\) are therefore built from
  \(\psi\)-classes, evaluation pullbacks, an equivariant Euler class and
  proper pushforward on smooth spaces, and they restrict to the corresponding
  classes of \(T_{N}\).  When the fixed locus is compact, the equivariant
  \(J\)-function defined through the localized pairing, which is the
  definition used in \cite{SS}, has this pushforward form in degree
  \(d\ge1\), because \(\mathrm{ev}\) factors through the zero section and
  \(j_{T}^{*}j_{T*}=e_{G}(V'(-1))\cup-\).  The extremal specialization of
  \(T_{N}\) kills the curve classes of \(B_{N}\).  The class \(c_{1}(B_{N})\)
  pairs to zero with \([L]\), so by the divisor equation the shift by
  \(-\ln(q)c_{1}(B_{N})\) multiplies both sides of \eqref{eq:ss35} by the same
  factor \(e^{-\ln(q)c_{1}(B_{N})/z}\).

  Proposition~\ref{prop:IJ}, which allows non-split bundles, now gives the
  identity in \(H^{*}(T_{N})\).  For fixed powers of \(q\), \(z\) and
  \(\ln q\) and a fixed monomial in the coordinates of \(t\), each side's
  coefficient is a finite sum of homogeneous components in
  \(H^{*}_{G}(T)\), and each component is detected in \(H^{*}(T_{N})\) once
  \(2N-1\) exceeds its degree.  Taking \(N\) large proves the identity there.
\end{proof}

The proof uses only the non-equivariant identity over arbitrary smooth
projective bases, so the same approximation gives the equivariant form of
\cite[Theorem~4.4]{SS} for \(r-s>1\) from its printed non-equivariant form.

\begin{remark}
  \label{rem:circularity}
  The route just taken differs from the printed attribution of
  \cite[Remark~4.5(3)]{SS} in one place.  The authors deduce cone membership
  from \cite{Brown}, \cite{CoatesGivental} and \cite[Lemma~9.6]{SS}.  The
  first two are the split-case inputs of Lemma~\ref{lem:cone} and carry no
  hypothesis on \(\nu\).  But \cite[Lemma~9.6]{SS} rests on
  \cite[Lemma~9.4]{SS}, whose identification of the fundamental solutions of
  \(T\) and its associated-graded model is deduced from \eqref{eq:ss35}
  read as a \(J\)-function formula, that is, from \cite[Theorem~4.4]{SS} with
  its hypothesis \(r-s>1\).  At \(\nu=1\) that formula is what is being
  established, so Lemma~\ref{lem:split-reduction} appeals to the underlying
  deformation instead of to the conclusion drawn from it.  Once
  Proposition~\ref{prop:IJ} and its \(S\)-equivariant form,
  Corollary~\ref{cor:equivariant}, are available, \cite[Lemmas~9.4--9.6]{SS}
  hold at \(\nu=1\) as printed, their input having been supplied; this is
  what Section~\ref{sec:scope} uses.
\end{remark}

\begin{corollary}
  \label{cor:presentation}
  Let \(r=s+1\) with \(s\ge1\).  Then the presentation \eqref{eq:ss37} of
  \cite[Theorem~4.6]{SS} holds, the modified extremal \(J\)-function
  \cite[(39)--(40)]{SS} of \(T\) is as printed, and
  \cite[Theorem~1.2]{SS} holds: after the extremal
  specialization, \(c_{1}(X)\star_{q}-\) on \(QH^{*}(X)\) has eigenvalue
  \(0\) with multiplicity \(\rk H^{*}(X')\) and the single nonzero
  eigenvalue
  \[
    \lambda_{0}(q)=e^{-\pi\sqrt{-1}\,s}q=(-1)^{s}q
  \]
  with multiplicity \(\rk H^{*}(Z)\).
\end{corollary}

\begin{proof}
  The proof of \cite[Theorem~4.6]{SS} specializes the parameter \(t\) of
  \cite[Theorem~4.4]{SS}, exhibits a differential operator annihilating
  \(J^{T}\), and concludes with \cite[Theorem~10.3.1]{CoxKatz} and
  \cite[Lemma~2]{Pandharipande}; it uses \(r-s>1\) only to know that
  \eqref{eq:ss35} is \(J^{T}\).  Proposition~\ref{prop:IJ} supplies that,
  so the proof applies verbatim, and \cite[(39)--(40)]{SS} follow by the
  same substitution as in \cite[Section~4.3]{SS}.  The same derivation,
  applied to Corollary~\ref{cor:equivariant}, gives the equivariant
  presentation of \(QH^{*}_{\mathrm{ext},\mathbb T}(T)\) stated after
  \cite[(81)]{SS}; that step is the authors' own, and only its input is
  supplied here.

  For the spectrum, \cite[Propositions~5.1 and~5.2]{SS} concern the abstract
  ring \eqref{eq:ss37} and its subspaces and assume only \(r>s\), and
  \cite[Corollary~5.3]{SS} combines them with the ring homomorphism
  \(i^{*}\colon QH^{*}_{\mathrm{ext}}(X)\to QH^{*}_{\mathrm{ext}}(T)\) of
  \cite[Theorem~3.3]{SS}.  Once \eqref{eq:ss37} is available at \(\nu=1\),
  that argument is unchanged.  Its output is \(r-s\) nonzero eigenvalues
  \(\lambda_{k}(q)=(r-s)e^{-\pi\sqrt{-1}(2k+s)/(r-s)}q\) for
  \(0\le k<r-s\); at \(r-s=1\) only \(k=0\) occurs, and
  \(\lambda_{0}(q)=e^{-\pi\sqrt{-1}s}q\).
\end{proof}

\begin{remark}
  \label{rem:pbundle}
  The companion statement \cite[(34)]{SS} for \(\PP(V)\) is the \(s=0\) case
  of the same expression, and Lemma~\ref{lem:order} applied to it gives
  maximal \(z\)-order \(1-rd\).  This is at most \(-1\) for every \(d\ge1\)
  as soon as \(r\ge2\), which holds whenever \(r=s+1\) with \(s\ge1\).  So
  \cite[(34), (36) and~(38)]{SS} are covered in the same range, by the same
  count.
\end{remark}

Proposition~\ref{prop:IJ} and Corollaries~\ref{cor:equivariant}
and~\ref{cor:presentation} prove Theorem~\ref{thm:normalization}.  They also settle the status of
\cite[Remark~4.5(3)]{SS}: under the standing assumption \(r>s\) its range
\(r-s\le1\) is exactly \(\nu=1\), and by
Corollary~\ref{cor:normalized} and Remark~\ref{rem:boundary} its negative
assertion holds there only for \(s=0\).
